% يجب الترجمة بمحرك LuaLaTeX  (Menu -> Compiler -> LuaLaTeX)
% الإجابة النموذجية لاختبار منتصف الفصل الأول - الصف السابع - الرياضيات
\documentclass[12pt, a4paper]{article}
\usepackage[a4paper, top=1.2cm, bottom=1.4cm, left=1.4cm, right=1.4cm]{geometry}
\usepackage{amsmath, amssymb}
\usepackage{fontspec}
\usepackage[bidi=basic]{babel}
\babelprovide[main, import, onchar=ids fonts]{arabic}
\babelprovide[import, onchar=ids fonts]{english}
\babelfont{rm}{Noto Serif}
\babelfont[arabic]{rm}{Amiri}
\usepackage{enumitem}
\usepackage{tabularx}
\usepackage{array}
\usepackage{graphicx}
\usepackage{tikz}
\usepackage{eso-pic}

\setlength{\parindent}{0pt}
\newcommand{\LTRfig}[1]{\mbox{\foreignlanguage{english}{#1}}}

% عنوان سؤال داخل شريط رمادي مع العلامة على اليسار
\newcommand{\question}[2]{%
  \par\vspace{0.2cm}\noindent
  \colorbox{gray!15}{\parbox{\dimexpr\linewidth-2\fboxsep\relax}{\textbf{\large #1}\hfill\mbox{\textbf{(#2)}}}}%
  \par\vspace{0.1cm}}
% علامة خطوة الحل
\newcommand{\m}[1]{\hfill\mbox{\textcolor{blue!60!black}{\textbf{[#1]}}}}
% الإجابة النهائية داخل إطار
\newcommand{\ans}[1]{\fbox{\textbf{#1}}}
% رأس الفرع
\newcommand{\sub}[1]{\par\vspace{0.12cm}\noindent\textbf{#1)}\ }

\begin{document}

% إطار الصفحات
\AddToShipoutPictureBG{%
  \LTRfig{%
  \begin{tikzpicture}[remember picture, overlay]
    \draw[line width=1.2pt] ([shift={(0.7cm,-0.7cm)}]current page.north west) rectangle ([shift={(-0.7cm,0.7cm)}]current page.south east);
  \end{tikzpicture}}%
}

% ============================ الترويسة ============================
\noindent
\begin{minipage}[t]{0.37\textwidth}
\vspace{0pt}
\textbf{المملكة الأردنية الهاشمية} \\
\textbf{وزارة التربية والتعليم} \\
\textbf{مدرسة عنبة الثانوية الشاملة للبنات} \\
\textbf{العام الدراسي: 2025 / 2026}
\end{minipage}%
\begin{minipage}[t]{0.26\textwidth}
\vspace{0pt}
\centering
\IfFileExists{logo.png}{\includegraphics[width=2.2cm]{logo.png}}{}
\end{minipage}%
\begin{minipage}[t]{0.37\textwidth}
\vspace{0pt}
\textbf{المبحث: الرياضيات} \\
\textbf{الصف: السابع} \\
\textbf{الزمن: حصة واحدة} \\
\textbf{الدرجة الكلية: 33 علامة}
\end{minipage}

\vspace{0.2cm}
\hrule height 1.2pt
\vspace{0.06cm}
\hrule height 0.4pt
\vspace{0.15cm}

\begin{center}
\textbf{\LARGE الإجابة النموذجية} \\[0.1cm]
\textbf{\large اختبار منتصف الفصل الأول}
\end{center}
\vspace{-0.2cm}
\noindent{\small الأرقام بين الأقواس المربعة \textcolor{blue!60!black}{\textbf{[ ]}} هي توزيع العلامة على خطوات الحل، والإجابة النهائية داخل إطار.}

% ============================ السؤال الأول ============================
\question{السؤال الأول: تمثيل الأعداد النسبية بالصورة العشرية والكسرية}{6 علامات}

\sub{أ} $\dfrac{3}{8} = 3 \div 8 = $ \ans{$0.375$} \quad (بالقسمة الطويلة) \m{1 للقسمة + 1 للناتج}

\sub{ب} $0.45 = \dfrac{45}{100} = \dfrac{45 \div 5}{100 \div 5} = $ \ans{$\dfrac{9}{20}$} \m{1 للكسر + 1 للتبسيط}

\sub{ج} $\dfrac{4}{9} = 4 \div 9 = 0.444\ldots = $ \ans{$0.\overline{4}$} \quad (الرقم 4 يتكرر) \m{1 للقسمة + 1 للناتج}

% ============================ السؤال الثاني ============================
\question{السؤال الثاني: ترتيب الأعداد النسبية ومقارنتها وجمعها}{6 علامات}

\sub{أ} نحوّل الأعداد إلى الصورة العشرية: \m{1}
\par\quad $-\dfrac{5}{6} \approx -0.833$ \ ، \ $-0.8$ \ ، \ $65\% = 0.65$ \ ، \ $\left|-\dfrac{3}{4}\right| = 0.75$
\par\quad الترتيب تصاعديًا: \ans{$-\dfrac{5}{6} \ ,\ -0.8 \ ,\ 65\% \ ,\ \left|-\dfrac{3}{4}\right|$} \m{1}

\sub{ب} نوحّد المقامات (المقام المشترك $56$): \quad
$-\dfrac{5}{8} = -\dfrac{35}{56}$ \ ، \ $-\dfrac{2}{7} = -\dfrac{16}{56}$ \m{1}
\par\quad بما أن $-35 < -16$ \ فإن \ \ans{$-\dfrac{5}{8} < -\dfrac{2}{7}$} \m{1}

\sub{ج} $-\dfrac{5}{6} + \dfrac{1}{4} = -\dfrac{10}{12} + \dfrac{3}{12}$ \m{1}
\par\quad $= $ \ans{$-\dfrac{7}{12}$} \m{1}

% ============================ السؤال الثالث ============================
\question{السؤال الثالث: ضرب الأعداد النسبية وقسمتها}{6 علامات}

\sub{أ} نحوّل العدد الكسري إلى كسر غير فعلي: \quad $-2\dfrac{1}{2} = -\dfrac{5}{2}$ \m{1}
\par\quad $-\dfrac{5}{2} \times \left(-\dfrac{4}{5}\right) = \dfrac{20}{10}$ \quad (سالب × سالب = موجب) \m{1}
\par\quad $= $ \ans{$2$} \m{1}

\sub{ب} $-2\dfrac{1}{4} = -\dfrac{9}{4}$ \m{1}
\par\quad $6 \div \left(-\dfrac{9}{4}\right) = 6 \times \left(-\dfrac{4}{9}\right) = -\dfrac{24}{9}$ \quad (نضرب في المقلوب) \m{1}
\par\quad $= $ \ans{$-\dfrac{8}{3} = -2\dfrac{2}{3}$} \m{1}

\newpage
% ============================ السؤال الرابع ============================
\question{السؤال الرابع: الأسس وترتيب العمليات الحسابية}{6 علامات}

\sub{أ} $\dfrac{1}{2} \times \dfrac{1}{2} \times 3 \times 3 \times 3 = $ \ans{$\left(\dfrac{1}{2}\right)^2 \times 3^3$}
\m{1 لكل أس صحيح}

\sub{ب} $4^{-2} = \dfrac{1}{4^2}$ \m{1}
\quad $= $ \ans{$\dfrac{1}{16}$} \m{1}

\sub{ج} $3(4-2)^3 + 8 \div 2^2 = 3(2)^3 + 8 \div 4$ \quad (الأقواس ثم الأسس) \m{1}
\par\quad $= 3 \times 8 + 2 = 24 + 2 = $ \ans{$26$} \quad (الضرب والقسمة ثم الجمع) \m{1}

% ============================ السؤال الخامس ============================
\question{السؤال الخامس: الجبر والتفكير العالي}{6 علامات}

\sub{أ} الحدود الجبرية: \ans{$5x \ ,\ -3y \ ,\ 8$} \m{1}
\par\quad معامل $y$: \ans{$-3$} \qquad الحد الثابت: \ans{$8$} \m{1}

\sub{ب} نبدأ من النهاية ونعكس كل عملية ($11\dfrac{4}{5} = 11.8$): \m{1}
\par\quad $8.9 + 11.8 = 20.7$ \quad ثم \quad $20.7 - 15 = 5.7$ \quad ثم \quad $5.7 + 6.3 = 12$
\par\quad الكمية الأصلية من الوقود: \ans{$12\ \mathrm{L}$} \m{1}
\par\quad {\small التحقق: \ $12 - 6.3 + 15 - 11.8 = 8.9$ \ (صحيح)}

\sub{ج} الترتيب: \ans{$\dfrac{1}{c} < \dfrac{1}{b} < \dfrac{1}{a}$} \m{1}
\par\quad مثال: \ $a = -4 \ ,\ b = -2 \ ,\ c = -1$ \ فإن \
$\dfrac{1}{a} = -\dfrac{1}{4} \ ,\ \dfrac{1}{b} = -\dfrac{1}{2} \ ,\ \dfrac{1}{c} = -1$
\par\quad و \ $-1 < -\dfrac{1}{2} < -\dfrac{1}{4}$ ، أي أن ترتيب المعكوسات الضربية للأعداد السالبة يكون عكس ترتيب الأعداد نفسها. \m{1}

% ============================ السؤال السادس ============================
\question{السؤال السادس: صح أم خطأ}{3 علامات}

{\small (نصف علامة لكل جملة)}
\vspace{0.05cm}

\renewcommand{\arraystretch}{1.25}
\noindent\begin{tabularx}{\linewidth}{|c|c|X|}
\hline
\textbf{الجملة} & \textbf{الإجابة} & \textbf{التوضيح} \\ \hline
1 & \textbf{T} & مثل: $0.37 = \dfrac{37}{100}$ \\ \hline
2 & \textbf{F} & كل عدد عشري دوري هو عدد نسبي، مثل: $0.\overline{3} = \dfrac{1}{3}$ \\ \hline
3 & \textbf{T} & مثل: $a = -1 \ ,\ b = -3$ \ فإن \ $|a| = 1 < |b| = 3$ \\ \hline
4 & \textbf{F} & المعكوس الضربي للعدد $-\dfrac{3}{4}$ هو $-\dfrac{4}{3}$ \\ \hline
5 & \textbf{F} & مثل: $-2 - (-5) = 3$ \ وهو عدد موجب \\ \hline
6 & \textbf{T} & $2^{-3} = \dfrac{1}{2^3} = \dfrac{1}{8}$ \\ \hline
\end{tabularx}

\vfill
\begin{center}
\textbf{مع تمنياتي لكم بالتوفيق}
\end{center}

\end{document}
